\chapter{Refinement Types and Bounds Check Elimination}
\label{chap:refinements}

\section{Path-Sensitive Numerical Refinement}
\label{sec:refinements:intro}

In high-performance systems programming languages, memory safety requires that every array indexing operation $A[i]$ is guarded by a bounds check verifying that $0 \le i < \text{len}(A)$. In hot numerical loops and stencil kernels, dynamic bounds checks incur branch misprediction penalties and inhibit vectorization.

NumLang resolves this overhead by embedding a path-sensitive \emph{Interval Refinement Domain} (\texttt{src/mir/supercompiler/state.rs}) directly into the symbolic driving engine. Whenever symbolic driving traverses a conditional branch or loop guard, the numerical bounds of all active places are refined along each path.

\section{The Interval Abstract Domain}
\label{sec:refinements:lattice}

We formalize the interval abstract domain over signed 64-bit integers:
\begin{equation}
\mathcal{I} = \{ [\ell, u] \mid \ell, u \in \mathbb{Z} \cup \{-\infty, +\infty\}, \ell \le u \} \cup \{ \bot \}
\end{equation}
The partial order $\sqsubseteq$ corresponds to subset inclusion:
\begin{equation}
[\ell_1, u_1] \sqsubseteq [\ell_2, u_2] \iff \ell_2 \le \ell_1 \land u_1 \le u_2
\end{equation}

\subsection{Lattice Operations}
The join ($\sqcup$) and meet ($\sqcap$) operators are defined as:
\begin{align}
[\ell_1, u_1] \sqcup [\ell_2, u_2] &= [\min(\ell_1, \ell_2), \max(u_1, u_2)] \\
[\ell_1, u_1] \sqcap [\ell_2, u_2] &= [\max(\ell_1, \ell_2), \min(u_1, u_2)] \quad (\text{yielding } \bot \text{ if } \max > \min)
\end{align}

\subsection{Widening Operator}
To guarantee termination across loop headers and knot generalizations, NumLang defines the widening operator $\nabla$:
\begin{equation}
[\ell_1, u_1] \nabla [\ell_2, u_2] = \left[
\begin{cases}
\ell_1 & \text{if } \ell_1 \le \ell_2 \\
-\infty & \text{otherwise}
\end{cases}, \quad
\begin{cases}
u_1 & \text{if } u_1 \ge u_2 \\
+\infty & \text{otherwise}
\end{cases}
\right]
\end{equation}

\section{Abstract Transfer Functions for Arithmetic Operations}
\label{sec:refinements:transfer}

The \texttt{Interval} domain implements abstract transfer functions for all MIR binary operations (Phase 33):

\begin{align}
[\ell_1, u_1] \oplus_+ [\ell_2, u_2] &= [\ell_1 + \ell_2, u_1 + u_2] \\[0.6em]
[\ell_1, u_1] \oplus_- [\ell_2, u_2] &= [\ell_1 - u_2, u_1 - \ell_2] \\[0.6em]
[\ell_1, u_1] \oplus_\times [\ell_2, u_2] &= [\min(P), \max(P)] \quad \text{where } P = \{\ell_1 \ell_2, \ell_1 u_2, u_1 \ell_2, u_1 u_2\} \\[0.6em]
[\ell_1, u_1] \oplus_{/\!} [c, c] &= \left[\left\lfloor \frac{\ell_1}{c} \right\rfloor, \left\lfloor \frac{u_1}{c} \right\rfloor\right] \quad (c > 0) \\[0.6em]
[\ell_1, u_1] \oplus_{\ll} [k, k] &= [\ell_1 \ll k, u_1 \ll k] \quad (k \ge 0, \ell_1 \ge 0) \\[0.6em]
[\ell_1, u_1] \oplus_{\gg} [k, k] &= [\ell_1 \gg k, u_1 \gg k] \quad (k \ge 0)
\end{align}

\section{Branch Narrowing and Dead-Branch Pruning}
\label{sec:refinements:narrowing}

When the driver encounters $\mathtt{BranchIf}(x < c, bb_{\text{then}}, bb_{\text{else}})$:
\begin{enumerate}
    \item \textbf{Then-Branch Refinement}: The environment updates $x \gets x \sqcap [-\infty, c - 1]$.
    \item \textbf{Else-Branch Refinement}: The environment updates $x \gets x \sqcap [c, +\infty]$.
    \item \textbf{Dead-Branch Pruning}: If $x \sqcap [-\infty, c - 1] = \bot$, the entire then-branch is pruned statically from the process tree. Conversely, if $x \sqcap [c, +\infty] = \bot$, the else-branch is pruned.
\end{enumerate}

\section{Static Bounds Check Elimination (BCE)}
\label{sec:refinements:bce}

When evaluating an array indexing instruction $x = A[i]$ where $A$ has length $L$:
\begin{enumerate}
    \item The driver queries the interval environment for $i$: let $\mathcal{I}(i) = [\ell_i, u_i]$.
    \item If $0 \le \ell_i$ and $u_i < L$, the condition $0 \le i < L$ is proven invariant.
    \item The compiler emits an unchecked memory load directly, completely omitting the panic check basic block.
    \item The global diagnostic counter \texttt{sc\_bce\_eliminated} is incremented.
\end{enumerate}

In the Supercompiler Showdown benchmark suite, this optimization eliminates 100\% of runtime bounds checks across all array-based algorithms (\texttt{kmp}, \texttt{fib\_matrix}, \texttt{map\_map}).

\section{Interval Propagation and Abstract Interpretation}
\label{sec:refinement:interval_impl}

The interval lattice implements standard abstract interpretation operators: meet ($\sqcap$), join ($\sqcup$), and widening ($\nabla$). Let $I_1 = [l_1, u_1]$ and $I_2 = [l_2, u_2]$. The lattice transfer functions over $\mathbb{Z}_\infty$ are defined as:
\begin{align}
I_1 \sqcap I_2 &= [\max(l_1, l_2), \min(u_1, u_2)] \\
I_1 \sqcup I_2 &= [\min(l_1, l_2), \max(u_1, u_2)] \\
I_1 \nabla I_2 &= \left[ \begin{cases} l_1 & \text{if } l_1 \le l_2 \\ -\infty & \text{otherwise} \end{cases}, \begin{cases} u_1 & \text{if } u_1 \ge u_2 \\ +\infty & \text{otherwise} \end{cases} \right]
\end{align}

When evaluating comparison branches such as $x < c$, the true branch narrows the interval of $x$ via $I_x \sqcap [-\infty, c - 1]$, while the false branch narrows via $I_x \sqcap [c, +\infty]$. If any refined interval becomes empty ($l > u$), the branch is provably unreachable and completely pruned from the residual control-flow graph. Array indexing expressions $A[i]$ are verified by querying whether $I_i \sqsubseteq [0, \mathrm{len}(A) - 1]$; if true, the hardware bounds check is safely omitted.
